\section*{Chapter \ref{ch:b1:hotspots} --- Hotspots and Intraplate Volcanism}

\begin{solution}{exo:b1:hotspots:1}
(a) $4205 + 5274 = \qty{9479}{m}$, nearly \qty{9.5}{km}. The volcano is
heavier than the water it replaces: the plate bends down beneath its load, so
its base lies deeper than the floor far away, and its true height is greater.
(b) $74000/0.1 = \num{740000}$ years at Kilauea's supply;
$74000/0.025 = \num{3.0e6}$ years at the mean rate of the chain. (c)
$0.12\times19 = \qty{2.3}{km^3}$, about $\num{3e-5}$ of Mauna Loa: three
thousandths of one per cent.
\end{solution}

\begin{solution}{exo:b1:hotspots:2}
(a) $0.92\times\pi/180 = \qty{0.0161}{rad.Myr^{-1}}$. (b) In radians per
million years times kilometres, $\omega R\sin\Delta$ is in kilometres per
million years, that is millimetres a year: $0.016057\times6371 = 102.3$, so
$v = 102.3\sin\Delta$, which gives \qty{51.1}{mm.yr^{-1}} at \ang{30},
\qty{88.6}{mm.yr^{-1}} at \ang{60}, \qty{102.3}{mm.yr^{-1}} at \ang{90} and
\qty{97.3}{mm.yr^{-1}} at \ang{108.0}. (c) The speed is proportional to the
distance $R\sin\Delta$ from the axis, largest on the \omterm{def:b1:spherical-astronomy:spherical-triangle}{great circle} \ang{90}
from the pole; beyond it the point comes closer to the axis again, on the
side of the opposite pole.
\end{solution}

\begin{solution}{exo:b1:hotspots:3}
(a) Deccan: $\num{e6}/\num{5e5} = \qty{2}{km}$. Central Atlantic:
$\num{2e6}/\num{7e6} = \qty{0.29}{km}$, about \qty{290}{m}. (b) The mean for
the Deccan is close to its \qty{2.5}{km} section: the province is a thick,
continuous plateau. (c) The Central Atlantic
area is that of the whole region, on four continents, over which its
remnants are found: lavas, dykes and sills now separated by an ocean and much
eroded. The rock is concentrated in some places and absent from most of the
area, so the mean says nothing of a single pile.
\end{solution}

\begin{solution}{exo:b1:hotspots:4}
(a) $\rho_m\alpha\,\Delta T = 3300\times\num{3.2e-5}\times250 =
\qty{26.4}{kg.m^{-3}}$. (b) The swell column has the extra mass
$(\rho_m - \rho_w)h = 2300\times690 = \qty{1.587e6}{kg.m^{-2}}$, which the
light layer must balance: $H = \num{1.587e6}/26.4 = \qty{6.0e4}{m}$, about
\qty{60}{km} of hot rock.
\end{solution}

\begin{solution}{exo:b1:hotspots:5}
(a) At A ($\lambda = 0$, $\varphi = 0$), $\vec e = (0, 1, 0)$ and
$\vec n = (0, 0, 1)$; with $\sin\ang{270} = -1$ and $\cos\ang{270} = 0$,
$\vec t_A = -\vec e = (0, -1, 0)$. At B ($\lambda = \ang{90}$),
$\vec e = (-1, 0, 0)$, $\vec n = (0, 0, 1)$ and $\vec t_B = (1, 0, 0)$.
(b) $\vec t_A\times\vec t_B = (0, 0, 1)$. (c) With $\vec p' = (0, 0, 1)$ and
$\vec r_A = (1, 0, 0)$, $\vec p'\times\vec r_A = (0, 1, 0)$, and
$(\vec p'\times\vec r_A)\cdot\vec t_A = -1 < 0$: the pole is
$-\vec p' = (0, 0, -1)$, the South Pole. A westward motion along the equator
is clockwise seen from above the North Pole, so anticlockwise seen from above
the South Pole. (d) At \ang{90} from the pole, $\omega = 50/6371 =
\qty{7.85e-3}{rad.Myr^{-1}} = \qty{0.45}{\degree\per\Myr}$.
\end{solution}

\begin{solution}{exo:b1:hotspots:6}
(a) $\sum x_i v_i = 90000 + 240000 + 342000 = \num{672000}$ and
$\sum x_i^2 = \num{7.0e7}$, so $\omega = \qty{0.0096}{rad.Myr^{-1}}$, or
\qty{0.55}{\degree\per\Myr}. (b) The predictions are 28.8, 48.0 and
\qty{57.6}{mm.yr^{-1}}; the residuals $+1.2$, $0.0$ and $-0.6$
\unit{mm.yr^{-1}}. (c) For a rigid rotation the speed is proportional to the
distance from the axis, and a point on the axis does not move: there is no
intercept to fit.
\end{solution}

\begin{solution}{exo:b1:hotspots:7}
(a) $v = \qty{50e-3}{m}/\qty{3.156e7}{s} = \qty{1.58e-9}{m.s^{-1}}$, so
$B = 2300\times\num{3e8}\times\num{1.58e-9} = \qty{1090}{kg.s^{-1}}$ and
$Q = (910/\num{3.2e-5})\times1090 = \qty{3.1e10}{W}$. (b) A third of the
Hawaiian flux. (c) $B$ grows with $A\,v$: the plume under the faster plate
must remake the same swell twice as fast, and its flux is twice as large.
\end{solution}

\begin{solution}{exo:b1:hotspots:8}
(a) $350/(16.1 - 10.3) = 350/5.8 = \qty{60}{km.Myr^{-1}}$, that is
\qty{60}{mm.yr^{-1}}. (b) The young track, \qty{29 \pm 5}{mm.yr^{-1}}, agrees
with the \qty{26 \pm 5}{mm.yr^{-1}} of the plate within the uncertainties; the
old part went more than twice as fast, $60 - 26 = 34$, almost seven standard
uncertainties of the plate speed. (c) Yes: N\,\ang{54}\,E is the \omterm{def:b1:spherical-astronomy:horizontal}{azimuth}
\ang{54}, opposite to \ang{235} within a degree. The volcanism moves over the
plate the opposite way to the plate over the \omterm{def:b1:hotspots:hotspot}{hotspot}. (d) The stretching moved
the older calderas apart, lengthening the distances between them and so the
apparent speed; the plate itself did not go faster.
\end{solution}

\begin{solution}{exo:b1:hotspots:9}
(a) $\num{9.2e10}/\num{e12} = \qty{0.092}{W.m^{-2}}$, \qty{92}{mW.m^{-2}}.
(b) Nine to eighteen times the observed excess, which may even be nil.
(c) $L^2/\kappa =
\num{e10}/\num{e-6} = \qty{e16}{s}$, about \qty{320}{Myr}; the plate crosses
\qty{1000}{km} in $1000/92 = \qty{11}{Myr}$. Heat from below has had time to
cross only a small part of the plate while it is over the swell: most of the
plume's heat stays in the \omterm{def:b1:layered-earth:layers}{mantle} beneath the plate and is carried away with
it, and the surface sees little of it.
\end{solution}

\begin{solution}{exo:b1:hotspots:10}
(a) Ten times the melt, \qty{e7}{km^3}, of radius
$(3V/4\pi)^{1/3} = \qty{134}{km}$. (b) $\num{e7}/100 = \qty{e5}{km^2}$, a disc
$2\sqrt{\num{e5}/\pi} = \qty{357}{km}$ across: eighteen times smaller in area
than the Traps. The head was larger or melted less, or it spread wider and
thinner, or the lavas flowed far beyond it. (c) \qty{1}{km^3} a year, ten
times the present supply of Kilauea. (d) A new plume must gather enough
hot, light rock to push its way up through the viscous \omterm{def:b1:layered-earth:layers}{mantle}, and a large
volume rises far faster than a small one; the tail that follows flows up the
hot, weak path the head has opened, and a narrow conduit suffices.
\end{solution}

\begin{solution}{exo:b1:hotspots:11}
(a) $\qty{40}{km.Myr^{-1}}\times(81 - 47)\,\unit{Myr} = \qty{1360}{km}$.
(b) The plate moves at $(-80\sin\ang{20}, 80\cos\ang{20}) = (-27.4, 75.2)$
and the \omterm{def:b1:hotspots:hotspot}{hotspot} at $(0, -40)$: relative to the \omterm{def:b1:hotspots:hotspot}{hotspot} the plate moves at
$(-27.4, 115.2)$ \unit{mm.yr^{-1}}, a speed of \qty{118.4}{mm.yr^{-1}} towards
$\ang{360} - \arctan(27.4/115.2) = \ang{346.6}$. (c) The track would give a
speed \qty{48}{\percent} too high and a direction \ang{6.6} too far east:
a drifting \omterm{def:b1:hotspots:hotspot}{hotspot} records the plate's motion relative to itself, not to the
deep \omterm{def:b1:layered-earth:layers}{mantle}.
\end{solution}

\begin{solution}{exo:b1:hotspots:12}
(a) $\qty{3e6}{J.m^{-3}.K^{-1}}\times\qty{1300}{K}\times\qty{e15}{m^3} =
\qty{3.9e24}{J}$. (b) Over $\qty{3.16e13}{s}$, \qty{1.2e11}{W}: about the
heat carried by the Hawaiian plume today (1.3 times), and \qty{0.26}{\percent}
of the \qty{47}{TW} of the Earth. (c) A lower bound: the head also brings up
the excess heat of the nine tenths of its rock that do not melt, and warms the
base of the \omterm{def:b1:layered-earth:lithosphere}{lithosphere}, and the latent heat was left out.
\end{solution}

\begin{solution}{pb:b1:hotspots:1}
\textbf{1.} At Hawaii ($\varphi = \ang{19.0}$, $\lambda = \ang{-155.2}$),
with $\sin\ang{304} = -0.8290$ and $\cos\ang{304} = 0.5592$:
\[
\begin{aligned}
  \vec r_H &= (-0.8583,\ -0.3966,\ 0.3256), &
  \vec e &= (0.4195,\ -0.9078,\ 0),\\
  \vec n &= (0.2955,\ 0.1366,\ 0.9455), &
  \vec t_H &= (-0.1825,\ 0.8289,\ 0.5287).
\end{aligned}
\]

\textbf{2.} At Louisville ($\varphi = \ang{-53.6}$, $\lambda = \ang{-140.6}$),
with $\sin\ang{316} = -0.6947$ and $\cos\ang{316} = 0.7193$:
\[
\begin{aligned}
  \vec r_L &= (-0.4586,\ -0.3767,\ -0.8049), &
  \vec e &= (0.6347,\ -0.7727,\ 0),\\
  \vec n &= (-0.6220,\ -0.5109,\ 0.5934), &
  \vec t_L &= (-0.8883,\ 0.1693,\ 0.4269).
\end{aligned}
\]

\textbf{3.} $\lVert\vec t_H\rVert^2 = 0.0333 + 0.6871 + 0.2795 = 0.9999$ and
$\vec t_H\cdot\vec r_H = 0.1566 - 0.3287 + 0.1721 = 0.0000$, to the rounding.

\textbf{4.} $\vec t_H\times\vec t_L = (0.2643, -0.3917, 0.7054)$, of length
$0.8491$. The angle between the tracks is $\arcsin0.8491 = \ang{58.1}$ (their
dot product, $0.528$, is positive, so the angle is acute).

\textbf{5.} $\vec p' = (0.3113, -0.4613, 0.8308)$. Then
$(\vec p'\times\vec r_H)\cdot\vec t_H = -0.98 < 0$: the pole is
$\vec p = (-0.3113, 0.4613, -0.8308)$.

\textbf{6.} $\arcsin(-0.8308) = \ang{-56.18}$ and
$\operatorname{atan2}(0.4613, -0.3113) = \ang{124.01}$: the pole lies at
\ang{56.2}\,S, \ang{124.0}\,E, in the Southern Ocean south of Australia.

\textbf{7.} $\vec p\cdot\vec r_H = 0.2672 - 0.1830 - 0.2705 = -0.1863$, so
$\Delta_H = \ang{100.7}$; $\vec p\cdot\vec r_L = 0.1428 - 0.1738 + 0.6687 =
0.6377$, so $\Delta_L = \ang{50.4}$.

\textbf{8.} The vector $\vec p\times\vec r_L$ is perpendicular to $\vec p$ and
to $\vec r_L$, as $\vec t_L$ is (by construction of $\vec p$, and because a
track is horizontal); both lie along the same line, and its length is
$\sin\Delta_L$. So $(\vec p\times\vec r_L)\cdot\vec t_L = \pm\sin\Delta_L$, and
the computation gives $+0.770 = \sin\ang{50.4}$: the plate moves at Louisville
along $+\vec t_L$, the way the track runs. The sign rule, applied at Hawaii,
could have given the wrong sense at Louisville; it did not, so both
directions are reproduced exactly, \ang{304} and \ang{316}.

\textbf{9.} $\omega = 67/(6371\sin\ang{50.38}) = 67/4908 =
\qty{0.01365}{rad.Myr^{-1}}$, that is \qty{0.782}{\degree\per\Myr}.

\textbf{10.} $360/0.782 = \qty{460}{Myr}$.

\textbf{11.} The cosine rule between \ang{56.18}\,S, \ang{124.01}\,E and
\ang{63.78}\,S, \ang{112.48}\,E gives \ang{9.5}. The two-track rate,
\qty{0.78}{\degree\per\Myr}, is \qty{15}{\percent} below the
\qty{0.92}{\degree\per\Myr} fitted to thirteen speeds.

\textbf{12.} Society \ang{73.6}, Pitcairn \ang{77.8}, Cobb \ang{134.7}.

\textbf{13.} $\omega R = 0.01365\times6371 = \qty{86.98}{mm.yr^{-1}}$, times
$\sin\Delta$: Society \qty{83.5}{mm.yr^{-1}}, Pitcairn
\qty{85.0}{mm.yr^{-1}}, Cobb \qty{61.8}{mm.yr^{-1}}.

\textbf{14.} Pitcairn agrees: \qty{5}{mm.yr^{-1}} less than measured, within
its \qty{15}{mm.yr^{-1}}. Society is \qty{25.5}{mm.yr^{-1}} too slow, 2.6
uncertainties, and Cobb \qty{18.8}{mm.yr^{-1}} too fast, more than six:
neither agrees.

\textbf{15.} The pole is perpendicular to both track vectors. If they are
nearly parallel, their cross product is short, and an error of a few degrees in
either direction turns it through a large angle: the pole is poorly
determined.

\textbf{16.} $v = \qty{92e-3}{m}/\qty{3.156e7}{s} = \qty{2.92e-9}{m.s^{-1}}$
and $B = 2300\times\num{4.83e8}\times\num{2.92e-9} = \qty{3240}{kg.s^{-1}}$.

\textbf{17.} $Q = (910/\num{3.2e-5})\times3240 = \qty{9.2e10}{W}$,
\qty{0.20}{\percent} of \qty{47}{TW}.

\textbf{18.} The ridges make $21/67000 = \qty{3.1e-4}{km^3}$ a year per
kilometre; $0.1/\num{3.1e-4} = \qty{320}{km}$ of ridge.

\textbf{19.} $74000/0.1 = \num{740000}$ years.

\textbf{20.} $\num{2e6}/\num{81e6} = \qty{0.025}{km^3}$ a year: Kilauea alone
is supplied today with four times the mean output of the whole \omterm{def:b1:hotspots:hotspot}{hotspot}, which
is near its strongest.

\textbf{21.} $v_H = 86.98\sin\ang{100.73} = 86.98\times0.9825 =
\qty{85.5}{mm.yr^{-1}}$.

\textbf{22.} $92 - 85.5 = \qty{6.5}{mm.yr^{-1}}$, 2.2 standard uncertainties of
the measured speed: the two \omterm{def:b1:hotspots:hotspot}{hotspots} do not quite agree on the motion of one
\omterm{def:b1:plate-kinematics:rigid-plate}{rigid plate} over a fixed frame.

\textbf{23.} The ages give \qty{9.1}{cm} and \qty{8.6}{cm} a year, averaged
over the whole chain; Morgan and Morgan's \qty{9.2}{cm} a year lies close to
the first, and the Louisville prediction, \qty{8.5}{cm} a year, close to the
second. All four lie within \qty{0.7}{cm} a year of one another: tracks, ages
and directions tell the same story to within about \qty{8}{\percent}.

\textbf{24.} To one decimal, the Louisville track gives the Pacific plate over
the Hawaiian \omterm{def:b1:hotspots:hotspot}{hotspot} the speed
\textbf{\qty[text-series-to-math=true]{8.5}{cm.yr^{-1}}}.
\end{solution}
